\documentclass[11pt]{article}
\usepackage[a4paper,margin=0.8in]{geometry}
\usepackage{calc}
\usepackage{array,longtable,booktabs}
\usepackage{xcolor}
\usepackage{hyperref}
\usepackage{enumitem}
\usepackage{fancyvrb}
\usepackage{fvextra}
\usepackage{framed}
\usepackage{color}
\fvset{breaklines=true, breakanywhere=true}
\DefineVerbatimEnvironment{Highlighting}{Verbatim}{breaklines,breakanywhere,commandchars=\\\{\}}
\usepackage{color}
\usepackage{fancyvrb}
\newcommand{\VerbBar}{|}
\newcommand{\VERB}{\Verb[commandchars=\\\{\}]}
\DefineVerbatimEnvironment{Highlighting}{Verbatim}{commandchars=\\\{\}}
% Add ',fontsize=\small' for more characters per line
\usepackage{framed}
\definecolor{shadecolor}{RGB}{248,248,248}
\newenvironment{Shaded}{\begin{snugshade}}{\end{snugshade}}
\newcommand{\AlertTok}[1]{\textcolor[rgb]{0.94,0.16,0.16}{#1}}
\newcommand{\AnnotationTok}[1]{\textcolor[rgb]{0.56,0.35,0.01}{\textbf{\textit{#1}}}}
\newcommand{\AttributeTok}[1]{\textcolor[rgb]{0.13,0.29,0.53}{#1}}
\newcommand{\BaseNTok}[1]{\textcolor[rgb]{0.00,0.00,0.81}{#1}}
\newcommand{\BuiltInTok}[1]{#1}
\newcommand{\CharTok}[1]{\textcolor[rgb]{0.31,0.60,0.02}{#1}}
\newcommand{\CommentTok}[1]{\textcolor[rgb]{0.56,0.35,0.01}{\textit{#1}}}
\newcommand{\CommentVarTok}[1]{\textcolor[rgb]{0.56,0.35,0.01}{\textbf{\textit{#1}}}}
\newcommand{\ConstantTok}[1]{\textcolor[rgb]{0.56,0.35,0.01}{#1}}
\newcommand{\ControlFlowTok}[1]{\textcolor[rgb]{0.13,0.29,0.53}{\textbf{#1}}}
\newcommand{\DataTypeTok}[1]{\textcolor[rgb]{0.13,0.29,0.53}{#1}}
\newcommand{\DecValTok}[1]{\textcolor[rgb]{0.00,0.00,0.81}{#1}}
\newcommand{\DocumentationTok}[1]{\textcolor[rgb]{0.56,0.35,0.01}{\textbf{\textit{#1}}}}
\newcommand{\ErrorTok}[1]{\textcolor[rgb]{0.64,0.00,0.00}{\textbf{#1}}}
\newcommand{\ExtensionTok}[1]{#1}
\newcommand{\FloatTok}[1]{\textcolor[rgb]{0.00,0.00,0.81}{#1}}
\newcommand{\FunctionTok}[1]{\textcolor[rgb]{0.13,0.29,0.53}{\textbf{#1}}}
\newcommand{\ImportTok}[1]{#1}
\newcommand{\InformationTok}[1]{\textcolor[rgb]{0.56,0.35,0.01}{\textbf{\textit{#1}}}}
\newcommand{\KeywordTok}[1]{\textcolor[rgb]{0.13,0.29,0.53}{\textbf{#1}}}
\newcommand{\NormalTok}[1]{#1}
\newcommand{\OperatorTok}[1]{\textcolor[rgb]{0.81,0.36,0.00}{\textbf{#1}}}
\newcommand{\OtherTok}[1]{\textcolor[rgb]{0.56,0.35,0.01}{#1}}
\newcommand{\PreprocessorTok}[1]{\textcolor[rgb]{0.56,0.35,0.01}{\textit{#1}}}
\newcommand{\RegionMarkerTok}[1]{#1}
\newcommand{\SpecialCharTok}[1]{\textcolor[rgb]{0.81,0.36,0.00}{\textbf{#1}}}
\newcommand{\SpecialStringTok}[1]{\textcolor[rgb]{0.31,0.60,0.02}{#1}}
\newcommand{\StringTok}[1]{\textcolor[rgb]{0.31,0.60,0.02}{#1}}
\newcommand{\VariableTok}[1]{\textcolor[rgb]{0.00,0.00,0.00}{#1}}
\newcommand{\VerbatimStringTok}[1]{\textcolor[rgb]{0.31,0.60,0.02}{#1}}
\newcommand{\WarningTok}[1]{\textcolor[rgb]{0.56,0.35,0.01}{\textbf{\textit{#1}}}}
\setlist[itemize]{leftmargin=*}
\setlist[enumerate]{leftmargin=*}
\hypersetup{colorlinks=true,linkcolor=blue,urlcolor=blue}
\providecommand{\tightlist}{%
  \setlength{\itemsep}{0pt}\setlength{\parskip}{0pt}}
\title{R Workshop -- Day 3 -- Exploratory Data Analysis and Statistical Graphics}
\author{Applied R for Statistical Teaching, Research and Reproducible Analysis}
\date{}
\begin{document}
\maketitle
\tableofcontents
\newpage
\hypertarget{r-workshop-day-3-exploratory-data-analysis-and-statistical-graphics}{%
\section{R Workshop -- Day 3 -- Exploratory Data Analysis and
Statistical
Graphics}\label{r-workshop-day-3-exploratory-data-analysis-and-statistical-graphics}}

\textbf{Session 3 of 7 \textbar{} Duration: 3 hours} \textbf{Programme:}
Applied R for Statistical Teaching, Research and Reproducible Analysis

\hypertarget{learning-objectives}{%
\subsection{Learning objectives}\label{learning-objectives}}

By the end of this session you will be able to:

\begin{enumerate}
\def\labelenumi{\arabic{enumi}.}
\tightlist
\item
  Start from a research question and choose a chart type that answers
  it, rather than choosing a chart type first.
\item
  Explain and use the grammar of graphics: data, aesthetics, geoms,
  scales, labels and themes.
\item
  Build and interpret bar charts, histograms, box plots and scatter
  plots in \texttt{ggplot2}.
\item
  Add titles, axis labels and units so a plot stands on its own.
\item
  Export publication-ready graphics to a file.
\item
  Write a visual claim that is accurate about what the chart can and
  cannot show.
\end{enumerate}

\hypertarget{capstone-link}{%
\subsection{Capstone link}\label{capstone-link}}

Today you produce the first visual evidence for your capstone report:
three exported, interpreted graphics built from the cleaned anchor
dataset (\texttt{Data/processed/anchor\_dataset\_clean.csv}, produced on
Day 2).

\hypertarget{segment-plan-3-hours}{%
\subsection{Segment plan (3 hours)}\label{segment-plan-3-hours}}

\begin{longtable}[]{@{}
  >{\raggedright\arraybackslash}p{(\columnwidth - 4\tabcolsep) * \real{0.3000}}
  >{\raggedleft\arraybackslash}p{(\columnwidth - 4\tabcolsep) * \real{0.4000}}
  >{\raggedright\arraybackslash}p{(\columnwidth - 4\tabcolsep) * \real{0.3000}}@{}}
\toprule\noalign{}
\begin{minipage}[b]{\linewidth}\raggedright
Segment
\end{minipage} & \begin{minipage}[b]{\linewidth}\raggedleft
Time
\end{minipage} & \begin{minipage}[b]{\linewidth}\raggedright
Purpose
\end{minipage} \\
\midrule\noalign{}
\endhead
\bottomrule\noalign{}
\endlastfoot
Opening and recap & 10 min & Reload the clean anchor dataset, recap the
research question \\
Concept bridge & 25 min & Question-led plotting and the grammar of
graphics \\
Guided coding & 60 min & Bar chart, histogram, box plot, scatter plot \\
Participant practice & 45 min & Build and label three graphics on the
anchor dataset \\
Interpretation and reporting & 25 min & Write a two- to three-sentence
interpretation per plot \\
Evidence log and capstone update & 15 min & Export and log the three
graphics \\
\end{longtable}

\begin{center}\rule{0.5\linewidth}{0.5pt}\end{center}

\hypertarget{part-1-concept-bridge-core-workshop-activity}{%
\subsection{Part 1 -- Concept bridge (Core workshop
activity)}\label{part-1-concept-bridge-core-workshop-activity}}

\hypertarget{start-with-the-question-not-the-chart}{%
\subsubsection{Start with the question, not the
chart}\label{start-with-the-question-not-the-chart}}

A chart type is a means, not an end. Three workshop-relevant questions
and their natural chart types:

\begin{longtable}[]{@{}
  >{\raggedright\arraybackslash}p{(\columnwidth - 2\tabcolsep) * \real{0.5000}}
  >{\raggedright\arraybackslash}p{(\columnwidth - 2\tabcolsep) * \real{0.5000}}@{}}
\toprule\noalign{}
\begin{minipage}[b]{\linewidth}\raggedright
Question
\end{minipage} & \begin{minipage}[b]{\linewidth}\raggedright
Natural chart
\end{minipage} \\
\midrule\noalign{}
\endhead
\bottomrule\noalign{}
\endlastfoot
How many lecturers are in each training track? & Bar chart of counts \\
How is post-training score distributed? & Histogram \\
Does post-training score differ by training track? & Box plot, score by
track \\
Is post-training score related to attendance? & Scatter plot \\
\end{longtable}

If you find yourself with a question and no obvious chart, it is worth
restating the question more precisely before opening \texttt{ggplot2}.

\hypertarget{the-grammar-of-graphics-briefly}{%
\subsubsection{The grammar of graphics,
briefly}\label{the-grammar-of-graphics-briefly}}

\texttt{ggplot2} builds a plot in layers:

\begin{itemize}
\tightlist
\item
  \textbf{data} -- the data frame.
\item
  \textbf{aesthetics (\texttt{aes()})} -- which columns map to which
  visual properties (x, y, colour, fill).
\item
  \textbf{geoms} -- the visual representation (\texttt{geom\_bar()},
  \texttt{geom\_histogram()}, \texttt{geom\_boxplot()},
  \texttt{geom\_point()}).
\item
  \textbf{scales} -- how data values map to visual ranges (axis breaks,
  colour palettes).
\item
  \textbf{labels and themes} -- titles, axis labels, units, and overall
  appearance.
\end{itemize}

Every \texttt{ggplot2} call you write this week follows the same shape:

\begin{Shaded}
\begin{Highlighting}[]
\FunctionTok{ggplot}\NormalTok{(data, }\FunctionTok{aes}\NormalTok{(}\AttributeTok{x =}\NormalTok{ ..., }\AttributeTok{y =}\NormalTok{ ...)) }\SpecialCharTok{+}
  \FunctionTok{geom\_...}\NormalTok{() }\SpecialCharTok{+}
  \FunctionTok{labs}\NormalTok{(}\AttributeTok{title =} \StringTok{"..."}\NormalTok{, }\AttributeTok{x =} \StringTok{"..."}\NormalTok{, }\AttributeTok{y =} \StringTok{"..."}\NormalTok{) }\SpecialCharTok{+}
  \FunctionTok{theme\_minimal}\NormalTok{()}
\end{Highlighting}
\end{Shaded}

\begin{center}\rule{0.5\linewidth}{0.5pt}\end{center}

\hypertarget{part-2-guided-coding-core-workshop-activity}{%
\subsection{Part 2 -- Guided coding (Core workshop
activity)}\label{part-2-guided-coding-core-workshop-activity}}

\begin{Shaded}
\begin{Highlighting}[]
\FunctionTok{library}\NormalTok{(tidyverse)}

\NormalTok{lecturers }\OtherTok{\textless{}{-}} \FunctionTok{read\_csv}\NormalTok{(}\StringTok{"Data/processed/anchor\_dataset\_clean.csv"}\NormalTok{, }\AttributeTok{show\_col\_types =} \ConstantTok{FALSE}\NormalTok{)}
\end{Highlighting}
\end{Shaded}

\hypertarget{bar-chart-counts-by-training-track}{%
\subsubsection{2.1 Bar chart -- counts by training
track}\label{bar-chart-counts-by-training-track}}

\begin{Shaded}
\begin{Highlighting}[]
\FunctionTok{ggplot}\NormalTok{(lecturers, }\FunctionTok{aes}\NormalTok{(}\AttributeTok{x =}\NormalTok{ training\_track)) }\SpecialCharTok{+}
  \FunctionTok{geom\_bar}\NormalTok{(}\AttributeTok{fill =} \StringTok{"\#2C7FB8"}\NormalTok{) }\SpecialCharTok{+}
  \FunctionTok{labs}\NormalTok{(}
    \AttributeTok{title =} \StringTok{"Participants by training track"}\NormalTok{,}
    \AttributeTok{x =} \StringTok{"Training track"}\NormalTok{, }\AttributeTok{y =} \StringTok{"Number of participants"}
\NormalTok{  ) }\SpecialCharTok{+}
  \FunctionTok{theme\_minimal}\NormalTok{()}
\end{Highlighting}
\end{Shaded}

\emph{Interpretation example:} ``Online was the most common track,
followed by In-Person and then Blended. This affects how much weight an
Online-driven result should carry in later analysis.''

\hypertarget{histogram-distribution-of-post-training-scores}{%
\subsubsection{2.2 Histogram -- distribution of post-training
scores}\label{histogram-distribution-of-post-training-scores}}

\begin{Shaded}
\begin{Highlighting}[]
\FunctionTok{ggplot}\NormalTok{(lecturers, }\FunctionTok{aes}\NormalTok{(}\AttributeTok{x =}\NormalTok{ post\_score)) }\SpecialCharTok{+}
  \FunctionTok{geom\_histogram}\NormalTok{(}\AttributeTok{binwidth =} \DecValTok{5}\NormalTok{, }\AttributeTok{fill =} \StringTok{"\#41AB5D"}\NormalTok{, }\AttributeTok{colour =} \StringTok{"white"}\NormalTok{) }\SpecialCharTok{+}
  \FunctionTok{labs}\NormalTok{(}
    \AttributeTok{title =} \StringTok{"Distribution of post{-}training scores"}\NormalTok{,}
    \AttributeTok{x =} \StringTok{"Post{-}training score (0{-}100)"}\NormalTok{, }\AttributeTok{y =} \StringTok{"Number of participants"}
\NormalTok{  ) }\SpecialCharTok{+}
  \FunctionTok{theme\_minimal}\NormalTok{()}
\end{Highlighting}
\end{Shaded}

\texttt{binwidth} controls how fine-grained the histogram is. Try 2, 5
and 10 and discuss how the same data can look smoother or noisier purely
from a plotting choice -- this is itself an honesty-in-graphics point.

\hypertarget{box-plot-score-by-training-track}{%
\subsubsection{2.3 Box plot -- score by training
track}\label{box-plot-score-by-training-track}}

\begin{Shaded}
\begin{Highlighting}[]
\FunctionTok{ggplot}\NormalTok{(lecturers, }\FunctionTok{aes}\NormalTok{(}\AttributeTok{x =}\NormalTok{ training\_track, }\AttributeTok{y =}\NormalTok{ post\_score)) }\SpecialCharTok{+}
  \FunctionTok{geom\_boxplot}\NormalTok{(}\AttributeTok{fill =} \StringTok{"\#FEB24C"}\NormalTok{) }\SpecialCharTok{+}
  \FunctionTok{labs}\NormalTok{(}
    \AttributeTok{title =} \StringTok{"Post{-}training score by training track"}\NormalTok{,}
    \AttributeTok{x =} \StringTok{"Training track"}\NormalTok{, }\AttributeTok{y =} \StringTok{"Post{-}training score (0{-}100)"}
\NormalTok{  ) }\SpecialCharTok{+}
  \FunctionTok{theme\_minimal}\NormalTok{()}
\end{Highlighting}
\end{Shaded}

A box plot shows the median, the interquartile range and likely outliers
in one view -- a natural bridge into the Day 5 question of whether track
is associated with score.

\hypertarget{scatter-plot-score-against-attendance}{%
\subsubsection{2.4 Scatter plot -- score against
attendance}\label{scatter-plot-score-against-attendance}}

\begin{Shaded}
\begin{Highlighting}[]
\FunctionTok{ggplot}\NormalTok{(lecturers, }\FunctionTok{aes}\NormalTok{(}\AttributeTok{x =}\NormalTok{ attendance, }\AttributeTok{y =}\NormalTok{ post\_score)) }\SpecialCharTok{+}
  \FunctionTok{geom\_point}\NormalTok{(}\AttributeTok{alpha =} \FloatTok{0.6}\NormalTok{, }\AttributeTok{colour =} \StringTok{"\#6A51A3"}\NormalTok{) }\SpecialCharTok{+}
  \FunctionTok{labs}\NormalTok{(}
    \AttributeTok{title =} \StringTok{"Post{-}training score by attendance"}\NormalTok{,}
    \AttributeTok{x =} \StringTok{"Attendance (\%)"}\NormalTok{, }\AttributeTok{y =} \StringTok{"Post{-}training score (0{-}100)"}
\NormalTok{  ) }\SpecialCharTok{+}
  \FunctionTok{theme\_minimal}\NormalTok{()}
\end{Highlighting}
\end{Shaded}

\hypertarget{export}{%
\subsubsection{2.5 Export}\label{export}}

\begin{Shaded}
\begin{Highlighting}[]
\FunctionTok{dir.create}\NormalTok{(}\StringTok{"Outputs"}\NormalTok{, }\AttributeTok{showWarnings =} \ConstantTok{FALSE}\NormalTok{)}

\NormalTok{p\_track\_counts }\OtherTok{\textless{}{-}} \FunctionTok{ggplot}\NormalTok{(lecturers, }\FunctionTok{aes}\NormalTok{(}\AttributeTok{x =}\NormalTok{ training\_track)) }\SpecialCharTok{+}
  \FunctionTok{geom\_bar}\NormalTok{(}\AttributeTok{fill =} \StringTok{"\#2C7FB8"}\NormalTok{) }\SpecialCharTok{+}
  \FunctionTok{labs}\NormalTok{(}\AttributeTok{title =} \StringTok{"Participants by training track"}\NormalTok{,}
       \AttributeTok{x =} \StringTok{"Training track"}\NormalTok{, }\AttributeTok{y =} \StringTok{"Number of participants"}\NormalTok{) }\SpecialCharTok{+}
  \FunctionTok{theme\_minimal}\NormalTok{()}

\FunctionTok{ggsave}\NormalTok{(}\StringTok{"Outputs/day3\_bar\_track\_counts.png"}\NormalTok{, p\_track\_counts, }\AttributeTok{width =} \DecValTok{7}\NormalTok{, }\AttributeTok{height =} \DecValTok{5}\NormalTok{, }\AttributeTok{dpi =} \DecValTok{300}\NormalTok{)}
\end{Highlighting}
\end{Shaded}

Always export at a fixed \texttt{width}/\texttt{height}/\texttt{dpi} so
figures look the same in your report regardless of your current RStudio
window size.

\hypertarget{writing-a-cautious-visual-claim}{%
\subsubsection{2.6 Writing a cautious visual
claim}\label{writing-a-cautious-visual-claim}}

A chart shows association, not cause. Compare:

\begin{itemize}
\tightlist
\item
  \emph{Overclaiming:} ``Attendance increases post-training scores.''
\item
  \emph{Accurate:} ``Participants with higher attendance tended to have
  higher post-training scores in this sample; this plot does not
  establish that attendance caused the difference.''
\end{itemize}

\begin{center}\rule{0.5\linewidth}{0.5pt}\end{center}

\hypertarget{part-3-participant-practice-core-workshop-activity}{%
\subsection{Part 3 -- Participant practice (Core workshop
activity)}\label{part-3-participant-practice-core-workshop-activity}}

Using \texttt{lecturers}, produce and export \textbf{three} graphics,
each with a two- to three-sentence interpretation:

\begin{enumerate}
\def\labelenumi{\arabic{enumi}.}
\tightlist
\item
  One categorical-summary chart (bar chart) for a variable of your
  choice (e.g.~\texttt{region}, \texttt{institution\_type},
  \texttt{completed}).
\item
  One distribution chart (histogram or box plot) for \texttt{pre\_score}
  or \texttt{post\_score}.
\item
  One relationship chart (scatter plot) connecting two numeric
  variables, or a box plot connecting a category to a numeric variable,
  relevant to your research question.
\end{enumerate}

For each: add a title, axis labels with units, export to
\texttt{Outputs/}, and write your interpretation as a comment directly
above the \texttt{ggsave()} call.

A facilitator-reviewed solution is in
\texttt{Solution-Scripts/day3\_eda\_solution.R}.

\begin{center}\rule{0.5\linewidth}{0.5pt}\end{center}

\hypertarget{optional-demonstration}{%
\subsection{Optional demonstration}\label{optional-demonstration}}

\begin{itemize}
\tightlist
\item
  \textbf{Density plots and violin plots} -- smoother alternatives to
  histograms/box plots; flag that both can mislead with very small
  groups (e.g.~fewer than \textasciitilde10 observations per group).
\item
  \textbf{Two-way categorical graphics} -- stacked or grouped bar
  charts, e.g.~\texttt{completed} by \texttt{training\_track}.
\item
  \textbf{Faceting} --
  \texttt{facet\_wrap(\textasciitilde{}\ training\_track)} to repeat one
  chart across groups.
\item
  \textbf{Missing-data visualisation} -- a simple bar chart of
  \texttt{colSums(is.na(lecturers))}.
\item
  \textbf{Honest versus misleading graphics} -- truncated y-axes, 3D pie
  charts, dual axes with mismatched scales.
\end{itemize}

\begin{Shaded}
\begin{Highlighting}[]
\CommentTok{\# Optional: faceted histogram of post\_score by training\_track}
\FunctionTok{ggplot}\NormalTok{(lecturers, }\FunctionTok{aes}\NormalTok{(}\AttributeTok{x =}\NormalTok{ post\_score)) }\SpecialCharTok{+}
  \FunctionTok{geom\_histogram}\NormalTok{(}\AttributeTok{binwidth =} \DecValTok{5}\NormalTok{, }\AttributeTok{fill =} \StringTok{"\#41AB5D"}\NormalTok{, }\AttributeTok{colour =} \StringTok{"white"}\NormalTok{) }\SpecialCharTok{+}
  \FunctionTok{facet\_wrap}\NormalTok{(}\SpecialCharTok{\textasciitilde{}}\NormalTok{ training\_track) }\SpecialCharTok{+}
  \FunctionTok{labs}\NormalTok{(}\AttributeTok{title =} \StringTok{"Post{-}training score by training track"}\NormalTok{, }\AttributeTok{x =} \StringTok{"Post{-}training score"}\NormalTok{, }\AttributeTok{y =} \StringTok{"Count"}\NormalTok{) }\SpecialCharTok{+}
  \FunctionTok{theme\_minimal}\NormalTok{()}
\end{Highlighting}
\end{Shaded}

\begin{center}\rule{0.5\linewidth}{0.5pt}\end{center}

\hypertarget{reference-or-take-home-material}{%
\subsection{Reference or take-home
material}\label{reference-or-take-home-material}}

\begin{itemize}
\tightlist
\item
  Theme customisation: \texttt{theme()}, \texttt{theme\_bw()},
  \texttt{theme\_classic()}, institutional colour palettes.
\item
  Annotation approaches: \texttt{geom\_text()}, \texttt{geom\_label()},
  \texttt{annotate()}.
\end{itemize}

\begin{center}\rule{0.5\linewidth}{0.5pt}\end{center}

\hypertarget{daily-deliverable}{%
\subsection{Daily deliverable}\label{daily-deliverable}}

Three exported graphics from the anchor dataset
(\texttt{Outputs/day3\_*.png}), each with a two- to three-sentence
interpretation saved alongside the script.

\hypertarget{evidence-log-entry}{%
\subsubsection{Evidence log entry}\label{evidence-log-entry}}

Complete the Day 3 row in \texttt{Workshop-Evidence-Log.md}: list the
three graphics produced and note one early pattern in the data relevant
to your research question.

\begin{center}\rule{0.5\linewidth}{0.5pt}\end{center}

\hypertarget{facilitator-notes}{%
\subsection{Facilitator notes}\label{facilitator-notes}}

\begin{itemize}
\tightlist
\item
  If a participant's research question does not map naturally onto
  bar/histogram/box plot/scatter, help them restate the question rather
  than forcing an awkward chart type.
\item
  Watch for binwidth and bar-width choices made purely for aesthetics
  without checking whether they hide or exaggerate a pattern -- this is
  a good moment to circle back to the ``honest versus misleading
  graphics'' discussion even if it is technically optional.
\item
  Confirm every exported PNG actually has a title and axis labels before
  participants move on; an unlabelled plot is treated as incomplete for
  the deliverable.
\end{itemize}
\end{document}
